\section*{Abstract}
\addcontentsline{toc}{section}{Abstract}

In Senegal, children remain deprived across a range of dimensions that undermine
their well-being, with more than half of them affected by multidimensional
poverty in 2018, while close to 24\% of households receive migrant remittances.
This paper estimates the causal impact of remittances on child multidimensional
poverty, using the two rounds of the Harmonized Household Living Standards
Survey (EHCVM~I, 2018-2019 and EHCVM~II, 2021-2022). Child poverty is measured
through UNICEF's multidimensional (MODA) approach, based on seven dimensions
broken down into 14 indicators, and three age groups are constructed (0-4, 5-14
and 15-17 years) to capture differentiated effects along the child's
development. Identification combines propensity score matching (PSM) and
difference-in-differences (PSM-DiD estimator of Heckman et al., 1997-1998),
jointly addressing selection on observables and time-invariant unobservable
heterogeneity. Treated households are those receiving, in 2018, a transfer from
a sender living outside Senegal who had previously been a member of the
household, which isolates migration originating from the household from other
forms of solidarity; they are compared with non-recipient households of
identical observable profile. Descriptively, the share of children deprived in
at least four dimensions fell from 56.6\% in 2018-2019 to 48.8\% in 2021-2022.
The average treatment effect on the treated is close to zero and not
significant for N-MODA incidence (ATT~=~-0.001; $p = 0.961$): over a three-year
horizon, deprivation among children of recipient households declined at the
same pace as among children of comparable non-recipient households, a
conclusion confirmed by the raw difference-in-differences and stable across
matching methods. This aggregate null result nonetheless masks two significant
effects of opposite sign: educational deprivation falls (-3.9~pp;
$p = 0.005$) while child protection deteriorates (5.1~pp; $p = 0.003$), the
five remaining dimensions being unchanged; since remittances are inseparable
from the departure of a household member, the monetary inflow invested in
schooling is offset by the costs of parental absence. The conclusion does not
depend on the convention used to define poverty: varying the cross-dimensional
threshold from $k=3$ to $k=6$ dimensions out of seven leaves the estimated
effect close to zero and non-significant at every threshold, while the analysis
of long-term recipients suggests that effects materialize beyond the three-year
horizon. Receiving remittances is therefore not sufficient, in the short run,
to alter the trajectory of children's multidimensional deprivations in Senegal:
these results call for lower remittance transaction costs, longer-term support
of recipient households so that flows translate into durable assets, stronger
complementarity between private transfers and public provision of basic
services, since structural deprivations cannot be lifted by purchasing power
alone, and targeted support for children whose parent is absent through
migration.

\bigskip

\noindent \textbf{Keywords}: child multidimensional poverty, N-MODA, remittances,
Senegal, EHCVM, propensity score matching (PSM), difference-in-differences (DiD).
